\section{Framework}
\label{sec:method}

A question $q$ comes with retrieved passages $P=(p_0,\dots,p_{n-1})$. A \emph{detector} returns flagged pairs
$F\subseteq\{(i,j): i<j\}$. A \emph{directional} detector also returns, for each flagged pair, a \emph{suspect}
set $S_{ij}\subseteq\{i,j\}$: the side it believes is misinformation. A \emph{format} turns $(P,F,S)$ into the
generator's prompt. We vary detectors and formats independently. Every prompt uses one skeleton: passages,
question, and an instruction to end with \texttt{Final answer(s): a | b | \dots}, listing every valid answer when
different entities share a name (Appendix~\ref{app:prompts}).

\subsection{Formats}
\label{sec:formats}
\begin{description}[leftmargin=0pt,itemsep=1pt,topsep=2pt]
\item[No flags (plain).] Passages and question only.
\item[Instruction.] Plain plus a generic instruction to check the passages for contradictions, with no detector.
We test two wordings. v1 ends ``if a disagreement reflects genuinely different valid answers, report each one'';
v2 asks the model to judge which passages are reliable and to leave out answers supported only by unreliable ones.
\item[Flag.] Each flagged passage gets the header \texttt{[POTENTIAL CONFLICT with passage(s): j]}, followed by a list
\texttt{Passages i and j: possible contradiction}. This is the format of the original CAAG system, except
that every detector's flags are rendered identically (no detector-specific type labels or confidence scores), so
detectors differ only in \emph{which} pairs they flag.
\item[Directional flag.] As Flag, but suspect passages are headed \texttt{[LIKELY MISINFORMATION; conflicts with
passage(s): j]}, and the list says which passage likely contains misinformation.
\item[Remove (suspect).] The suspect passages of every flagged pair are deleted from the prompt. For a pair without a
suspect, the second passage is deleted, which is what the original CAAG \textsc{remove} strategy does when there
are no retrieval scores.
\end{description}

\subsection{Detectors}
\label{sec:detectors}
\begin{description}[leftmargin=0pt,itemsep=1pt,topsep=2pt]
\item[CAAG] (Contradiction-Aware Attention Gating) is a heuristic detector from an earlier, unpublished RAG system, used here with its code unchanged. A pair is scored only if its MiniLM
\citep{reimers2019sbert} cosine similarity reaches $\theta_s=0.65$. Its score is the sum of a negation and contrast
asymmetry term, a named-entity mismatch term (differing dates, numbers, money or percentages) and a temporal-marker
term, multiplied by the similarity. A pair is flagged if the score reaches $\theta_c=0.35$. It runs locally and
gives no direction.
\item[NLI] is \texttt{nli-deberta-v3-base} \citep{he2023debertav3}, a cross-encoder scored in both directions; a pair is
flagged if the larger contradiction probability exceeds $0.5$. It gives no direction.
\item[LLM] is a single gpt-4o-mini call (temperature 0) that lists contradicting pairs, telling them apart from
ambiguity, complementary and temporal relations, and names the side that is more likely misinformation (or
``both'' or ``unclear''), as JSON.
\item[Random] flags each pair independently at CAAG's per-pair flag rate on the confirmatory split
($6.22\%$). It is a placebo that matches CAAG's flag volume.
\item[Oracle] flags exactly the gold contradiction pairs and names the gold misinformation side.
\end{description}

\subsection{Synthetic detectors}
\label{sec:synthetic}
To separate a detector's accuracy from its identity, we build detectors with a chosen pair-level precision $p$,
recall $r$ and direction accuracy $d$ from the gold labels. Each gold contradiction pair is flagged with probability $r$.
False flags are drawn from non-contradicting pairs at a single per-pair rate, set over the whole split so that
expected precision is $p$. That rate is dataset-wide, so conflict-free questions receive false flags too. A true flag
names the gold misinformation side with probability $d$, and otherwise the other side. A false flag has no
misinformation side, so its suspect is picked uniformly. Each pair gets one uniform draw for flagging and one for
direction, and every cell reuses them. Cells are therefore nested: raising $r$ or $d$, or lowering $p$, only adds
flags or corrections. This keeps comparisons across cells paired.

\subsection{Placing real detectors on the grid}
\label{sec:placement}
Can the grid predict a real detector's effect from a few summary numbers? We compare two ways of summarising a
detector. \emph{Pair-level}: its pair precision, recall and direction accuracy. \emph{Document-level}: what it does
to each subset's documents in a given format. For conflict questions, that is the share whose misinformation documents
are all marked and the share that have a correct document marked; for clean questions, the share that have a correct
document marked. Here ``marked'' means removed (Remove) or headed as misinformation (Directional flag). The prediction
for a subset is the inverse-distance-weighted mean effect of the four nearest grid cells in that space.
